\subsection{Editorial verification: what the root-multiplicity question actually asks}

No item in records 530--537 is assigned class C or D. The following proofs and computations are annotations checking the reconstruction and the literal errors; they are not claimed as new research. In particular, an unspecified wish for an illuminating multiplicity formula is not turned into an open demand for a positive formula. The symmetrizable formula in \citet[Theorem 2]{bermanmoody1979multiplicities} was already available in 1979. The elementary argument below covers the finite computational interpretation of 1981-533.4 without importing that paper's symmetrizability hypothesis.

\begin{proposition}[Recursive lowering determines every specified positive degree]\ours
Let \(A\) be a finite generalized Cartan matrix over characteristic zero. Use a Cartan realization with independent simple roots, and define the contragredient algebra by factoring the Chevalley cross-relation algebra by its largest ideal disjoint from the Cartan subalgebra. If \(\alpha\in Q_+\) has height greater than one, then
\[
 \mathfrak g_\alpha\longrightarrow
 \bigoplus_{i:\,\alpha-\alpha_i\in Q_+\setminus\{0\}}
 \mathfrak g_{\alpha-\alpha_i},\qquad
 v\longmapsto([f_i,v])_i
\]
is injective. Consequently a prescribed root-space dimension and independent representing words are computable by finitely many exact linear operations, inductively on height.
\end{proposition}
\begin{proof}
The positive part of the algebra before the maximal quotient is the free Lie algebra on the \(e_i\); its component of degree \(\alpha\) is finite-dimensional. The triangular decomposition and maximal-ideal construction are those of \citet[Chapter II, \S1]{kac1968graded}. The simple positive components have dimension one.

Suppose \(v\in\mathfrak g_\alpha\) is killed by all \(f_i\). Let \(J\) be the span of all iterated brackets
\(
 [e_{i_1},[e_{i_2},\ldots,[e_{i_s},v]\ldots]]
\), including \(v\). It is stable under the \(e_i\) and the Cartan subalgebra. It is stable under each \(f_j\) as well: move \(\operatorname{ad}f_j\) through the iterated positive adjoints, using
\([f_j,e_i]=-\delta_{ij}h_i\), and then move each Cartan adjoint through the remaining positive adjoints. Every resulting term is a scalar multiple of a shorter positive iterated bracket on \(v\); the terminal \([f_j,v]\) vanishes. This induction also covers an arbitrary linear combination of such brackets. Hence \(J\) is an ideal of \(\mathfrak g\).

Every weight of \(J\) has degree \(\alpha+\beta\) for \(\beta\in Q_+\), so \(J\cap\mathfrak h=0\). The maximal-quotient definition prohibits any nonzero such ideal: its preimage would strictly enlarge the largest ideal disjoint from \(\mathfrak h\). Therefore \(v=0\), proving injectivity. Terms with a negative coefficient in \(\alpha-\alpha_i\) are automatically zero; height greater than one excludes a Cartan-valued target. Removing the redundant central directions of Rudakov's enlarged realization does not change these positive components.

For effectiveness, choose a finite basis or spanning list of free Lie words of degree \(\alpha\). By induction suppose quotient coordinate maps are known in every smaller positive degree. The brackets \([f_i,w]\) are calculated from the cross relations and expressed in those already computed quotient coordinates. Stack the resulting matrices. Its kernel is exactly the degree-\(\alpha\) part of the maximal ideal, by the injection just proved. Its rank is the desired dimension, and a set of pivot columns supplies independent word representatives. All coefficients are obtained from the integer entries of \(A\), so exact rational arithmetic suffices for a generalized Cartan matrix. The calculation at every fixed height terminates because all free Lie components in question are finite-dimensional. No complexity bound, canonical choice of pivots or cancellation-free formula is asserted.
\end{proof}

\begin{remark}[Independent check of the lowering argument]\oursComputed
The supplied \texttt{editorial\_checks.py} computes the same quotient dimensions by two different procedures: the lowering-map recursion above, and a degree-by-degree Lie ideal generated by the two explicit Serre relations. Both use an exact free-Lie basis, but the second procedure never uses a lowering operator. For the symmetrizable rank-two pairs \((a,b)=(1,1),(2,1),(3,1),(2,2),(3,2)\), they agree through height eight. The first three height profiles are respectively
\[
 (2,1,0,0,0,0,0,0),\quad
 (2,1,1,0,0,0,0,0),\quad
 (2,1,1,1,1,0,0,0).
\]
These independently recover total dimensions \(3,4,6\). The affine pair \((2,2)\) gives \((2,1,2,1,2,1,2,1)\); the indefinite pair \((3,2)\) gives \((2,1,2,2,3,3,6,7)\). Finite output does not prove infinite-dimensionality; that assertion uses the previously established Cartan classification. The second procedure is used only for symmetrizable rank-two matrices, where the Serre presentation is justified. It is not substituted for the maximal-ideal definition in the general statement above.
\end{remark}

\subsection{Editorial counterexamples for the transformation-group notes}

For 1981-534.1--5 write \(G_{(S)}\) for the pointwise stabilizer and \(G_{\{S\}}\) for the setwise stabilizer. Then the source's closure is
\(
 \operatorname{cl}(S)=\operatorname{Fix}(G_{(S)}).
\)
This distinction is also needed for the restriction homomorphisms in the notebook's fourth axiom. The exact match of those axioms to Seifert's 1984 theorem remains unverified; the following corrections concern elementary statements, independent of that publication-identification gap.

\begin{proposition}[Singleton closure need not mean equal stabilizers]\ours
For every group action,
\[
 \operatorname{cl}\{x\}=\{y:G_x\subseteq G_y\}.
\]
Even in a faithful transitive action the containment can be strict. If the transitive set is finite, it is necessarily equality.
\end{proposition}
\begin{proof}
The displayed equivalence is just the definition that every element fixing \(x\) also fixes \(y\). For strictness let \(G=\mathbb Z[1/2]\rtimes\langle t\rangle\), with \(tbt^{-1}=2b\) on the additive translation subgroup, let \(H=\mathbb Z\), and let \(G\) act on left cosets \(G/H\). At \(x=H\), the stabilizer is \(H\). At \(y=t^{-1}H\), it is \(t^{-1}Ht=\tfrac12\mathbb Z\), strictly larger. Therefore \(y\in\operatorname{cl}\{x\}\) while \(G_y\ne G_x\). The action is transitive by construction and faithful: its kernel lies in \(\bigcap_{n\ge0}2^n\mathbb Z=\{0\}\), hence is trivial. Thus faithfulness does not repair the notebook's statement.

In a finite transitive set, \(G_x\) and \(G_y\) have the same finite index in \(G\). If \(G_x\subseteq G_y\), the index multiplication formula implies \([G_y:G_x]=1\). This is a separate verification of the missing finite-action hypothesis, and does not require \(G\) itself to be finite.
\end{proof}

\begin{example}[The degree-two exception to the double-transitivity remark]\ours
Let \(C_2\) act regularly on its two elements. It is transitive on the two ordered pairs of distinct elements, hence doubly transitive. Every point stabilizer is trivial, so \(\operatorname{cl}\{x\}=X\), strictly larger than \(\{x\}\). Thus the notebook's proposed obstruction requires a restriction. For a doubly transitive set with at least three points, \(G_x\) acts transitively on \(X\setminus\{x\}\), a set of at least two elements, and no point there is fixed by all of \(G_x\). In that range \(\operatorname{cl}\{x\}=\{x\}\) does hold.
\end{example}

\subsection{Independent finite checks of combinatorial formulas}

\begin{proposition}[Interval inversion]\ours
For interval multiplicities \(m_{pq}\) of an equioriented type-A representation,
\[
 r_{ij}=\sum_{p\le i\le j\le q}m_{pq},\qquad
 m_{ij}=r_{ij}-r_{i-1,j}-r_{i,j+1}+r_{i-1,j+1},
\]
where ranks beyond the ends are zero.
\end{proposition}
\begin{proof}
An interval summand contributes one to the rank of the composition from vertex \(i\) to vertex \(j\) exactly when it contains both endpoints, proving the first formula. The difference \(r_{ij}-r_{i-1,j}\) retains exactly intervals starting at \(i\) and ending at or beyond \(j\). Subtracting the corresponding difference with endpoint \(j+1\) retains exactly the intervals ending at \(j\). This proves the second formula. It verifies the indexing in the source completion from \citet[\S2.3]{abeasisdelfrakraft1981geometry}, not the deeper converse degeneration theorem.
\end{proof}

\begin{remark}[Independent enumeration results]\oursComputed
The code constructs explicit direct-sum arrow matrices for all 120 choices of three distinct intervals on four vertices, with independently varied multiplicities one or two. Matrix products and exact ranks agree with the finite-difference reconstruction in every case. For the free Lie formula, direct enumeration of primitive binary necklaces agrees with the M\"obius formula in all 65 bidegrees of total degree one through ten. The independent binary-tree recurrence agrees with the Catalan specialization for \(n=0,\ldots,12\).

For the Gelfand--Tsetlin question, recursive enumeration of every intervening integer row agrees with the Weyl product for all 210 weakly increasing top rows drawn from \(\{-2,-1,0,1,2,3,4\}\). Examples give counts \(1\) for \((0,0,0,0)\), \(4\) for \((0,0,0,1)\), \(64\) for \((0,1,2,3)\), and \(630\) for the additional row \((-2,-1,1,4)\). The general theorem is supplied by the interlacing basis \citep[\S3]{gelfandtsetlin1950unimodular} and the Weyl dimension formula; bounded enumeration is an independent check on indexing, endpoints and normalization.

Finally, multiplying a symmetric-algebra Euler product and separately expanding its formal logarithm through degree fourteen confirms \(C_n=\sum_{d\mid n}d\dim V_d\). This checks that the notebook's rescaled coefficients \(c_k=k\dim V_k\), including the denominator \(k\) in \(p(t)\), are internally consistent. Complete parameter ranges and output are retained beside the code.
\end{remark}

\subsection{Two sign and scaling checks}

\begin{example}[A direct test of the BGG rho shift]\ours
With \([e,f]=h\), \([h,f]=-2f\), let \(M(0)\) be the \(\mathfrak{sl}_2\) Verma module generated by \(v_0\), with \(ev_0=hv_0=0\). The nonzero vector \(fv_0\) has weight \(-2\), and \(e(fv_0)=hv_0=0\). The universal property gives a nonzero map \(M(-2)\to M(0)\); on the polynomial PBW bases it is multiplication by \(f\), so it is injective. Since \(\rho(h)=1\), the weights satisfy \(-2+1=-(0+1)\). They do not satisfy \(-2-1=\pm(0-1)\). This verifies the plus-rho convention against both an explicit module calculation and the shifted definition in \citet[\S3]{bernsteingelfandgelfand1971structure}.
\end{example}

\begin{proposition}[Determinant weights of the supergravity density]\ours
In the determinant-monomial ansatz \(A=|W|^p|\Gamma|^q\), the defining-equation and odd-frame covariance laws used by \citet[\S3, equation (8)]{schwarz1982supergravity} force \(p=4/3\), \(q=-1/3\).
\end{proposition}
\begin{proof}
Replacing the four real defining equations by \(f'=\eta f\) multiplies each determinant by \(\det\eta\), while the product of delta functions has weight \(-1\). Thus \(p+q=1\). Changing the complex odd tangent frame by a \(2\times2\) matrix \(S\) gives \(|W|\) weight \(|\det S|\) and the paired-spinor determinant \(|\Gamma|\) weight \(|\det S|^4\); hence \(p+4q=0\). The two equations have the unique solution stated.

As a separate check, rescale all defining equations by a positive real constant \(c\). Both four-by-four determinants acquire \(c^4\), and the delta product acquires \(c^{-4}\). The final density has total exponent \(16/3-4/3-4=0\), while the preliminary expression \(|\det\Gamma|^{-1/2}\prod\delta(f^k)\) has exponent \(-2-4=-6\). This confirms the source's internal correction. The argument determines exponents within this ansatz; it is not a replacement for the stronger local uniqueness classification attributed to \citet{gaydukromanovschwarz1981democracy}, whose full proof was not accessible here.
\end{proof}

\subsection{Verification of the Jantzen emendations}

This subsection checks two annotations to record 540. It does not claim a new solution of a C/D problem: the unrestricted direction question has the published negative answer in \citet[\S8.4]{williamson2016local}. The checks below provide independently reproducible evidence that the hypotheses and grading in the companion have not merely been copied from an abstract.

\begin{example}[Why moving lifts matter]\ours
Let \(A=\mathbb Q[t]_{(t)}\), and give \(A^2\) the symmetric bilinear form with matrix
\[
 S(t)=\begin{pmatrix}1&t\\t&t^2+t^3\end{pmatrix}.
\]
Its elementary divisors have orders \(0,3\). The class of \(e_2\) in \(A^2/tA^2\) therefore belongs to the third Jantzen-type step. The constant lift \(e_2\) has pairing \((t,t^2+t^3)\), whose order is only one. The moving lift \(e_2-te_1\) instead has pairing \((0,t^3)\), as required. This supplies an explicit reason why the literal fixed-vector definition on P000997 is incomplete. Direct multiplication and the determinant \(\det S(t)=t^3\) independently check the calculation.
\end{example}

\begin{example}[Exact reproduction of the type-A3 direction counterexample]\oursComputed
The file \texttt{verify\_annotations.py} builds the contravariant form from the matrix-unit relations
\[
 [E_{ij},E_{kl}]=\delta_{jk}E_{il}-\delta_{li}E_{kj}
\]
on a PBW basis of the weight space of highest weight \(-\alpha_1-\alpha_3\), lowered by \(2\alpha_1+3\alpha_2+2\alpha_3\). There are thirteen basis monomials. All arithmetic is exact. Dynkin coordinates of the highest weight are \((-2,2,-2)\); the three deformation directions are displayed below.
\[
\begin{array}{c|c|c}
\gamma(h_1),\gamma(h_2),\gamma(h_3)&
(\dim \operatorname{gr}^0,\dim \operatorname{gr}^1,
 \dim \operatorname{gr}^2,\dim \operatorname{gr}^3)&
(N(1),N(2),N(3),N(4))\\\hline
(1,1,1)&(7,3,2,1)&(6,9,10,10)\\
(1,-2,3)&(7,2,4,0)&(6,10,10,10)\\
(3,-2,1)&(7,2,4,0)&(6,10,10,10)
\end{array}
\]
Here \(N(k)\) is the nullity, over \(\mathbb Q\), of multiplication by the form matrix modulo \(t^k\). Every positive-coroot value of each direction is nonzero. The two non-dominant directions satisfy \(\gamma(h_1+2h_2+h_3)=0\); the third row independently checks the reflected choice.

Two different exact procedures agree. First, elimination over rational power series modulo \(t^5\) gives elementary-divisor orders. Second, fraction-free integer row elimination gives the ranks of the block matrices modulo \(t,t^2,t^3,t^4\). If the elementary-divisor orders are \(a_1,\ldots,a_{13}\), then
\[
 N(k)=\sum_{r=1}^{13}\min(k,a_r),\qquad
 \dim\operatorname{gr}^j=\#\{r:a_r=j\}.
\]
Stabilization between \(k=3\) and \(k=4\) also rules out hidden higher orders. The saved \texttt{verification\_output.json} and \texttt{verification\_run.log} record the full results. They reproduce the published counterexample and do not constitute a priority claim.
\end{example}

The same executable independently checks the signs in the Cauchy determinant at sizes 1 through 5, including a repeated-row edge case; the differentiated Rayleigh equation and its inverse transformation; an explicit nondegenerate binary-cubic Kirillov pairing; and all 511 finite surreal sign sequences of length at most eight by both canonical dyadic cuts and the corrected halving formula. These finite checks support particular emendations. They are not presented as proofs of an unbounded numerical conjecture.

\subsection{Haefliger's modular-group coinvariant equation: a complete continuous result}
\label{dev:1981-haefliger}

The question in 1981-555.1 is exact smooth solvability on the projective circle
\(X=\mathbb {RP}^{1}=\mathbb R\cup\{\infty\}\): for every real
\(f\in C^\infty(X)\), do there exist real \(u,v\in C^\infty(X)\) such that
\[
 f=u-u\circ T+v-v\circ S,
 \qquad T(x)=x+1,\qquad S(x)=-1/x?
\]
Both maps extend smoothly across the apparent affine-coordinate singularities.
The notebook calls exact solvability unknown and asserts density of the image
in the \(C^\infty\) topology. Those are two separate assertions. In particular,
density in a Fr\'echet space does not imply surjectivity.

\paragraph{Literature and the scope of this development.}
The search documented in the research log located Haefliger's 1980 paper on
minimal leaves, and the later invariant-functional framework of
\citet[abstract]{abouqatebelkacimi2000invariant}. The latter's
compact-group result does not cover this noncompact discrete group action.
Neither lead supplied a verified theorem answering the displayed smooth
equation. The subject connection with minimal leaves is recorded in the notes;
no historical explanation of the question's subsequent neglect is inferred.
The following result is a bounded contribution proved here, with no claim of
historical priority. It solves the exact equation on continuous functions. It
does not prove a regularity theorem for its solutions.

\begin{lemma}[No invariant probability measure]\ours
There is no Borel probability measure on \(X\) invariant under both \(T\)
and \(S\).
\end{lemma}
\begin{proof}
If \(\nu\) is \(T\)-invariant, the pairwise disjoint Borel sets
\([k,k+1)\), \(k\in\mathbb Z\), have the same measure. Finiteness forces
each measure to be zero, and countable additivity then gives
\(\nu(\mathbb R)=0\). Thus \(\nu=\delta_\infty\). Since
\(S(\infty)=0\), this probability is not \(S\)-invariant.
\end{proof}

\begin{lemma}[A uniform defect of invariance]\ours
Let \(\mathcal M(X)\) be the real finite signed regular Borel measures, with
total-variation norm \(\|\cdot\|_{\rm TV}\), and let \(g_*\mu\) denote
pushforward. There is a number \(\epsilon>0\) such that every
\(\mu\in\mathcal M(X)\) satisfies
\[
 \|\mu-T_*\mu\|_{\rm TV}+\|\mu-S_*\mu\|_{\rm TV}
 \ \geq\ \epsilon\|\mu\|_{\rm TV}.
 \tag{H1}
\]
\end{lemma}
\begin{proof}
The probabilities on compact \(X\) form a weak-* compact subset of
\(C(X)^*\). Pushforward by a homeomorphism is weak-* continuous and the
total-variation norm is weak-* lower semicontinuous. Consequently
\[
 \epsilon:=\inf_{\nu\ \mathrm{probability}}
 \bigl(\|\nu-T_*\nu\|_{\rm TV}+\|\nu-S_*\nu\|_{\rm TV}\bigr)
\]
is attained. If it were zero, a minimizer would contradict the preceding
lemma. Thus \(\epsilon>0\).

For signed measures \(\alpha,\beta\),
\(\||\alpha|-|\beta|\|_{\rm TV}\leq\|\alpha-\beta\|_{\rm TV}\).
Indeed, taking Radon--Nikodym derivatives with respect to
\(|\alpha|+|\beta|\) reduces this assertion to
\(\bigl||a|-|b|\bigr|\leq|a-b|\) and integration. A homeomorphism also
satisfies \(|g_*\mu|=g_*|\mu|\). Apply the positive-measure estimate to
\(|\mu|/\|\mu\|_{\rm TV}\), when \(\mu\ne0\), and use these two facts
for \(g=T,S\). The zero measure causes no exception.
\end{proof}

\begin{theorem}[Exact solvability for continuous functions]\ours
For every real \(f\in C(X)\), there are real \(u,v\in C(X)\) with
\[
 f=u-u\circ T+v-v\circ S.
 \tag{H2}
\]
If \(0<\epsilon'<\epsilon\), the solutions can be chosen to satisfy
\[
 \max(\|u\|_\infty,\|v\|_\infty)
 \leq \frac{2}{\epsilon'}\|f\|_\infty.
 \tag{H3}
\]
Here \(\epsilon\) is the positive constant in (H1); the theorem does not
claim an explicit numerical value for it or a linear choice of solutions.
\end{theorem}
\begin{proof}
Equip \(E=C(X)\oplus C(X)\) with the maximum norm and define the bounded
operator \(D:E\to C(X)\) by the right side of (H2). Under the Riesz
identification \(C(X)^*=\mathcal M(X)\), its adjoint is
\[
 D^*\mu=(\mu-T_*\mu,\mu-S_*\mu).
\]
The dual norm on \(E^*\) is the sum of the two total-variation norms.
Thus (H1) states precisely that \(D^*\) is bounded below by \(\epsilon\).

Let \(B_E\) be the closed unit ball and let
\(K=\overline{D(B_E)}\), with norm closure. For every \(\mu\),
\[
 \sup_{h\in K}\langle\mu,h\rangle
 =\|D^*\mu\|\geq\epsilon\|\mu\|_{\rm TV}.
\]
The Hahn--Banach separation theorem for closed convex sets implies
\(\epsilon' B_{C(X)}\subset K\) for every \(0<\epsilon'<\epsilon\).
To see the strict-inequality point explicitly, a hypothetical separating
functional for an element of \(\epsilon' B_{C(X)}\setminus K\) would
have supremum on \(K\) at least \(\epsilon\|\mu\|\), while its value
on that element is at most \(\epsilon'\|\mu\|\), a contradiction.

Set \(r_0=f\). Given \(r_j\ne0\), the inclusion just proved permits
a pair \(z_j\in E\) with
\[
 \|z_j\|_E\leq\|r_j\|_\infty/\epsilon',\qquad
 \|r_j-Dz_j\|_\infty\leq\tfrac12\|r_j\|_\infty.
\]
Put \(r_{j+1}=r_j-Dz_j\); if a residual is zero, take all subsequent
pairs zero. Then \(\|r_j\|_\infty\leq2^{-j}\|f\|_\infty\) and
\(\sum_j\|z_j\|_E\leq2\|f\|_\infty/\epsilon'\). Completeness gives
\(z=\sum_j z_j=(u,v)\in E\), and continuity of \(D\) gives
\(Dz=f\). This also proves (H3).
\end{proof}

\begin{proposition}[An exact reduction of the smooth question]\ours
For \(C_gw=w\circ g\), put \(A=(I+C_S)/2\). A smooth function \(f\)
has a smooth decomposition (H2) if and only if
\[
 Af\in\operatorname{ran}\bigl(A(I-C_T):C^\infty(X)\to
 C^\infty(X)^S\bigr),
 \tag{H4}
\]
where \(C^\infty(X)^S\) denotes the \(S\)-invariant functions.
\end{proposition}
\begin{proof}
Since \(S^2=I\) on \(\mathbb {RP}^{1}\), \(A\) is a projection and
\(A(I-C_S)=0\). Applying it to (H2) proves necessity. Conversely, choose
\(u\) with \(A(I-C_T)u=Af\). The remainder
\(r=f-(I-C_T)u\) satisfies \(Ar=0\), hence \(C_Sr=-r\).
Putting \(v=r/2\) gives \((I-C_S)v=r\) and proves sufficiency.
\end{proof}

\paragraph{Second check and remaining obstruction.}
An independent audit re-derived the measure argument and checked all norms.
It also gave a separate functional-analytic check: (H1) makes
\(\operatorname{ran}D^*\) norm closed and \(\ker D^*=0\); the Banach
closed-range theorem makes \(\operatorname{ran}D\) closed, while the
annihilator identity makes it dense, so \(D\) is onto. This agrees with
the constructive Hahn--Banach argument above. The use of total variation
of a signed measure is essential: mere absence of invariant signed
measures would show density, not a quantitative lower bound or
surjectivity.

The proof controls only uniform norms. Uniformly convergent series of
continuous solutions need not have bounded derivatives, and smoothing
does not commute with the two nonlinear coordinate actions in a way that
this argument estimates. In the smooth setting the dual obstructions are
distributions, not only measures. The notebook's claim of smooth density,
even if independently established, would still leave a closed-range or
regularity estimate to prove. Thus the strongest verified result in this
development is exact \(C^0\) solvability; the strongest formulation of the
original problem remains exact \(C^\infty\) solvability, retained as an
unresolved question rather than promoted to a theorem.

One further bounded investigation can target (H4): establish a tame
estimate for the averaged parabolic difference operator near the fixed
point \(\infty\), with a finite and explicitly controlled derivative
loss, or construct an obstruction to such an estimate. That is a precise
analytic task. The present argument supplies neither outcome, and the
absence of a verified resolution in this search is not evidence that no
resolution exists.

\subsection{Independent finite checks of two source emendations}

These checks support annotations 1981-548.1--2 and 1981-546.15; they are not
additional open problems. The partition theorem itself is the known identity
in \citet[Chapter I, Section 1, (1.7)--(1.7')]{macdonald1995symmetric}.

\begin{example}[Exact partition enumeration]\oursComputed
The supplied Python program enumerates partitions in every rectangle
\(0\leq m,n\leq6\) in two ways: weakly decreasing row lengths, and
\(n\)-element subsets of \(\{0,\ldots,m+n-1\}\), inverted as beta numbers.
Both enumerations agree and contain \(\binom{m+n}{n}\) partitions for each
rectangle. For all 3431 partitions across these 49 rectangles, the two
complementary exponent sets are disjoint and fill that interval. A separate
run on \(0\leq m,n\leq4\), with 251 partitions, agrees. The negative control
\(m=n=1\), \(\lambda=(1)\), makes the written generating-function left side
\(t+t^2\), whereas the required right side is \(1+t\). The corrected inverse
argument and conjugation give \(1+t\). These finite checks detect the source
error; they are not the proof for arbitrary rectangles.
\end{example}

\begin{proposition}[The simplex edge-length Jacobian]\ours
Order the independent Gram entries by diagonals first and then off-diagonals.
Write \(s_{ij}=\ell_{ij}^2\) for all \(n(n+1)/2\) edges of a labelled
\(n\)-simplex, with vertex zero distinguished. The absolute Jacobians are
\[
 \left|\det\frac{\partial(g_{ij})_{1\leq i\leq j\leq n}}
                         {\partial(s_{ij})_{0\leq i<j\leq n}}\right|
 =2^{-n(n-1)/2},
 \qquad
 \left|\det\frac{\partial(g_{ij})}{\partial(\ell_{ij})}\right|
 =2^n\prod_{0\leq i<j\leq n}\ell_{ij}.
\]
Consequently the density in equation (23) needs the product of lengths if
its differentials are length differentials.
\end{proposition}
\begin{proof}
The diagonal formula is \(g_{ii}=s_{0i}\). Each off-diagonal coordinate is
\(g_{ij}=(s_{0i}+s_{0j}-s_{ij})/2\). In the stated ordering the Jacobian is
block triangular, with diagonal blocks the identity and
\(-\tfrac12 I_{n(n-1)/2}\). This proves the first identity. The change from
lengths to squared lengths has diagonal Jacobian entries \(2\ell_{ij}\).
Their product gives the second identity. Finally
\(\det g=(n!\operatorname{vol}\sigma)^2\), so multiplying the invariant
matrix density by this Jacobian gives, up to a fixed constant,
\(\operatorname{vol}(\sigma)^{-(n+1)}\prod\ell_{ij}\,d\ell_{ij}\).
\end{proof}

\begin{example}[Second check of the Jacobian]\oursComputed
Independent symbolic differentiation of the full coordinate map in dimensions
\(n=1,2,3,4,5\) gives signed squared-length Jacobians
\(1,-1/2,-1/8,1/64,1/1024\), respectively. Division of the absolute
length-coordinate determinant by the product of all edge lengths gives
\(2,4,8,16,32\). An independent positive-definite test uses based orthonormal
edge vectors, whose pairwise squared distances are one or two, and recovers
the identity Gram matrix exactly in every dimension. For a single edge,
\(g=\ell^2\) gives \(g^{-1}dg=2\,d\ell/\ell\), while the written formula
would give \(d\ell/\ell^2\). The complete code and its output are included.
\end{example}

\subsection{Goldsheid's coordinate-product approximation: a counterexample and quantitative bounds}
\label{dev:1981-goldsheid}

\paragraph{The exact question and its two quantifiers.}
The intervention on P001018 (1981-557.8; GP-1981-40) concerns the standard
Euclidean symplectic space \(\mathbb R^{2n}\), with an orthonormal symplectic
basis \(e_1,\ldots,e_n,f_1,\ldots,f_n\). For every Lagrangian \(n\)-plane
\(L\), can one choose a nonzero vector in each prescribed plane
\(E_j=\operatorname{span}(e_j,f_j)\), so that their span \(K\) satisfies
\[
 \|P_L-P_K\|_{2\to2}\le 1/\sqrt2? \tag{G1}
\]
The notebook separately says that any constant \(c<1\), independent of
\(n\), would suffice. All projections here are Euclidean orthogonal
projections. Changing to a different, unquantified notion of distance would
change the problem.

\paragraph{Literature and historical limits.}
The searches in the accompanying log found no exact answer under the
speaker's name or the equivalent matrix formulation below. This is a
bounded negative retrieval result, not a proof that the problem has been
absent from the literature. A precise related result is
\citet[Theorem 3.4]{mehrmannpoloni2012doubling}: after symplectic coordinate
permutations, a Lagrangian has a graph matrix whose diagonal entries have
absolute value at most one and whose off-diagonal entries have absolute
value at most \(\sqrt2\). The theorem controls entries, and hence gives
dimension-dependent norm estimates; it does not state (G1) or a uniform
operator-gap estimate. We found no source explaining why Goldsheid's
specific question was dropped or whether it was later reformulated.

The results below are proved in this task, without a claim of historical
priority. They disprove (G1), show that any universal constant must be at
least \(\cos(\pi/8)\), and give an upper bound for every fixed dimension.
The still plausible weaker formulation is the existence of some
\(c\in[\cos(\pi/8),1)\) working in every dimension; no particular value
in that interval is asserted to work.

\begin{lemma}[Symmetric-unitary formulation]\ours
Identify \(\mathbb R^{2n}=\mathbb C^n\), with \(f_j=i e_j\), and let
\(Cz=\overline z\). Every real Lagrangian \(L\) is \(U\mathbb R^n\)
for a unitary matrix \(U\). Set \(S=UU^T\). Then \(S\) is symmetric
unitary and, as real-linear operators,
\[
 P_L=\tfrac12(I+SC).
\]
For \(K=D_\theta\mathbb R^n\), where
\(D_\theta=\operatorname{diag}(e^{i\theta_j})\),
\[
 \|P_L-P_K\|=\tfrac12\|S-D_\theta^2\|. \tag{G2}
\]
Consequently (G1) is equivalent to choosing a diagonal unitary \(D\)
such that \(\operatorname{Re}(S^*D)\) is positive semidefinite, where
\(\operatorname{Re}B=(B+B^*)/2\).
\end{lemma}
\begin{proof}
Choose a real orthonormal basis of \(L\). Its Hermitian Gram matrix is
the identity: the real part is the Euclidean Gram matrix and the imaginary
part vanishes because \(L\) is isotropic. Thus its columns form a unitary
\(U\). Reflection across \(L\) is the real-linear involution
\(UCU^*=(UU^T)C\). Averaging it with the identity gives the projection.
The same formula for \(D_\theta\) gives (G2), because conjugation is an
isometry and the norm of a complex-linear matrix is the same whether
regarded as a real or complex operator. Every diagonal unitary has a
diagonal unitary square root, so no choices are lost. Finally
\[
 (S-D)^*(S-D)=2I-S^*D-D^*S
             =2I-2\operatorname{Re}(S^*D),
\]
which proves the last equivalence.
\end{proof}

\begin{theorem}[An explicit six-plane refutes the proposed constant]\ours
Let \(\omega=e^{2\pi i/3}\) and
\[
 V=\frac1{\sqrt3}
 \begin{pmatrix}1&1&1\\1&\omega&\omega^2\\1&\omega^2&\omega\end{pmatrix}.
\]
The real six-dimensional subspace
\[
 L=\{(V\overline z,z):z\in\mathbb C^3\}\subset\mathbb C^6
       =\mathbb R^{12}
\]
is Lagrangian, and every coordinate-product Lagrangian \(K\) satisfies
\(\|P_L-P_K\|>1/\sqrt2\).
\end{theorem}
\begin{proof}
The Fourier matrix \(V\) is unitary. For \(z,w\in\mathbb C^3\), the
Hermitian inner product of the corresponding vectors in \(L\) is
\(\overline{z^*w}+z^*w\), which is real. The parametrization is real
linear and injective, so \(L\) is isotropic of half the ambient real
dimension and hence Lagrangian. An explicit real orthonormal basis is
\[
 \frac{(Ve_j,e_j)}{\sqrt2},\qquad
 \frac{(-iVe_j,ie_j)}{\sqrt2},\qquad j=1,2,3.
\]
Its reflection matrix is
\[
 S=\begin{pmatrix}0&V\\V^T&0\end{pmatrix}.
\]
Indeed \(S\) is symmetric unitary, and \(SC\) fixes exactly the displayed
plane: \(V^T\overline V=I\) supplies the second block equation.

Suppose a diagonal unitary \(D=\operatorname{diag}(d_1,\ldots,d_6)\)
made \(\|S-D\|\le\sqrt2\). The preceding lemma says
\(H=(S^*D+D^*S)/2\ge0\). Since every diagonal entry of \(S\) is zero,
every diagonal entry of \(H\) is zero. A positive semidefinite Hermitian
matrix obeys \(|H_{ij}|^2\le H_{ii}H_{jj}\), so \(H=0\).
For a nonzero entry \(S_{ij}\), this implies
\[
 d_i d_j=-\frac{S_{ij}}{\overline{S_{ij}}}. \tag{G3}
\]
The first row and first column of \(V\) are positive real. Their
equations force \(d_1=d_2=d_3=a\), \(d_4=d_5=d_6=b\), and \(ab=-1\).
The entry \(V_{22}=\omega/\sqrt3\) instead forces
\(ab=-\omega^2\), a contradiction. Thus the inequality cannot hold.
Compactness of the diagonal-unitary torus also shows that the minimum
gap for this one plane is strictly larger than \(1/\sqrt2\), not merely
that the infimum fails to be attained.
\end{proof}

\begin{theorem}[A Fourier-family obstruction to universal constants]\ours
For every odd integer \(m\ge3\), put
\(V_{jk}=m^{-1/2}e^{2\pi i jk/m}\), \(0\le j,k<m\), and form the
Lagrangian \(L_m=\{(V\overline z,z):z\in\mathbb C^m\}\subset\mathbb C^{2m}\).
Every coordinate-product Lagrangian \(K\) satisfies
\[
 \|P_{L_m}-P_K\|\ \geq\
 \frac12\sqrt{\,2+\sqrt{2-2/\sqrt m}\,}. \tag{G4}
\]
In particular any dimension-independent constant in Goldsheid's weaker
question must satisfy \(c\ge\cos(\pi/8)\).
\end{theorem}
\begin{proof}
As above, \(S=\bigl(\begin{smallmatrix}0&V\\V^T&0\end{smallmatrix}\bigr)\).
Write a diagonal unitary as \(D=\operatorname{diag}(A,B)\), with diagonal
entries \(a_j,b_k\). Then
\[
 \operatorname{Re}(S^*D)=\frac12
 \begin{pmatrix}0&C_0\\C_0^*&0\end{pmatrix},\qquad
 C_0=\overline V B+A^*V.
\]
The eigenvalues of the displayed off-diagonal Hermitian matrix are the
positive and negative singular values of \(C_0\), divided by two.
It follows exactly, without an estimate at this step, that
\[
 \|S-D\|^2=2+\|C_0\|. \tag{G5}
\]
Let \(W=(V_{jk}^2)\), the entrywise square, not the matrix square.
Because \(m\) is odd, multiplication by two permutes the residue classes
modulo \(m\). Consequently \(W\) is \(m^{-1/2}\) times a unitary
Fourier matrix with permuted columns, so \(\|W\|=m^{-1/2}\).
Expanding the Frobenius norm gives
\[
 \|C_0\|_F^2
 =2m+2\operatorname{Re}\sum_{j,k}\overline{a_j}\overline{b_k}V_{jk}^2
 \ge2m-2\sqrt m,
\]
because \(\|a\|=\|b\|=\sqrt m\) and
\(|a^*W\overline b|\le\|a\|\|W\|\|b\|=\sqrt m\).
There are \(m\) singular values, so
\(\|C_0\|\ge\|C_0\|_F/\sqrt m\ge\sqrt{2-2/\sqrt m}\).
Equations (G2) and (G5) prove (G4). Its right side tends to
\(\tfrac12\sqrt{2+\sqrt2}=\cos(\pi/8)\) as odd \(m\to\infty\).
Every \(c\) strictly below that limit is therefore defeated by some
finite odd \(m\). The endpoint itself is not excluded by this estimate.
\end{proof}

\begin{proposition}[A fixed-dimension upper bound]\ours
Every Lagrangian \(L\subset\mathbb C^n\) admits a common-phase product
plane \(K=e^{i\theta}\mathbb R^n\) with
\[
 \|P_L-P_K\|\le\cos\!\left(\frac{\pi}{2n}\right). \tag{G6}
\]
In particular the original bound is true for \(n=1,2\).
\end{proposition}
\begin{proof}
The eigenvalues of the unitary \(S\) lie on the unit circle. Of the \(n\)
circular gaps, including the wraparound gap and allowing repetitions,
one has length at least \(2\pi/n\). All eigenvalues therefore lie in a
closed complementary arc of length at most \(2\pi-2\pi/n\).
Let \(z=e^{i\phi}\) be its midpoint. Their angular distances from \(z\)
are at most \(\pi-\pi/n\), so, by the spectral theorem,
\[
 \|S-zI\|=\max_{\lambda\in\operatorname{Spec}S}|\lambda-z|
 \le2\sin\!\left(\frac\pi2-\frac{\pi}{2n}\right).
\]
Choose \(\theta=\phi/2\) and use (G2). If \(n=1\), the complementary
arc is a point and the bound is zero, as required.
\end{proof}

\begin{example}[Sharpness in dimension two]\ours
For \(S=\bigl(\begin{smallmatrix}0&1\\1&0\end{smallmatrix}\bigr)\),
each column of \(S-D\) has norm \(\sqrt2\) for every diagonal unitary
\(D\), so every product plane has gap at least \(1/\sqrt2\).
Taking \(D=iI\) gives equality, since \(S\) is self-adjoint unitary.
The associated Lagrangian is spanned over \(\mathbb R\) by
\((e_1+e_2)/\sqrt2\) and \((f_1-f_2)/\sqrt2\).
Thus (G6) is sharp at \(n=2\), although (G1) fails in higher dimension.
\end{example}

\begin{proposition}[The notebook's coordinatewise observation]\ours
For every coordinate plane \(E_j\), there is a unit vector
\(\varphi_j\in E_j\) with \(\operatorname{dist}(\varphi_j,L)\le1/\sqrt2\).
This does not bound the gap of the span of the chosen vectors.
\end{proposition}
\begin{proof}
Since \(L\) is Lagrangian, \(JL=L^\perp\), hence
\(P_L+JP_LJ^{-1}=I\). Compressing to the \(J\)-invariant plane \(E_j\)
and taking traces shows that the compression of \(P_L\) has trace one.
Its largest eigenvalue is at least \(1/2\), which proves the first
assertion by choosing a unit eigenvector.

For the second assertion, take \(n\ge2\),
\(v=(1,\ldots,1)/\sqrt n\in\mathbb R^n\), and
\(L=\mathbb Rv\oplus i v^\perp\). Its Hermitian pairings have zero
imaginary part and its real dimension is \(n\), so it is Lagrangian.
Every \(f_j=ie_j\) is at distance \(1/\sqrt n\) from \(L\), but
\(iv\in i\mathbb R^n\cap L^\perp\). Thus the span
\(K=i\mathbb R^n\) has \(\|P_L-P_K\|=1\). The equality follows also
from the fact that differences of orthogonal projections have norm at
most one. This example refutes an inference from a particular selection
of vectors, whereas the Fourier example refutes the existential (G1).
\end{proof}

\begin{remark}[Independent checks and numerical scope]\oursComputed
The independent mathematical audit re-derived the antiunitary reflection,
the positive-semidefinite zero-diagonal argument, the Fourier cross-term
estimate, the arc bound including \(n=1\), and the coordinatewise example.
It confirmed every hypothesis and emphasized that the endpoint
\(\cos(\pi/8)\) is not ruled out. A second representation of each Fourier
plane by the explicit real orthonormal basis checked isotropy directly.

The included \texttt{verify\_root.py} uses two independent random seeds,
five dimensions and four samples per dimension to compare the real
projection matrices with (G2); the maximum discrepancy was below
\(7\times10^{-16}\). It also checks Fourier examples for
\(m=3,5,7,11\). The rigorous lower bound for \(m=3\) is approximately
\(0.854312835963\), already above \(1/\sqrt2\). Two exploratory
optimizations found feasible gaps near \(0.866026\); neither met its
convergence criterion, and neither is reported as a global optimum.
The proofs of (G4)--(G6) do not depend on floating-point evidence. The
earlier random exploration, whose sampled gaps were all below
\(1/\sqrt2\), is also included: it illustrates why that evidence could
not justify the proposed universal bound.
\end{remark}

\paragraph{What remains and a bounded next investigation.}
Define
\[
 c_*:=\sup_{n\ge1}\ \sup_{S=S^T,\ S^*S=I}
       \inf_{D\ \mathrm{diagonal\ unitary}}\frac12\|S-D\|.
\]
The proofs establish \(\cos(\pi/8)\le c_*\le1\), with a strictly
subunit bound in each fixed dimension. Goldsheid's weaker question is
exactly whether \(c_*<1\). The original stronger assertion
\(c_*\le1/\sqrt2\) is false. A further bounded research run could
determine the exact optimum of the explicit \(m=3\) Fourier example by
an algebraic or interval certificate, or improve the Fourier-family
lower bound using higher singular-value moments. Either target is
concrete. A numerical optimizer alone would not decide the supremum over
all dimensions, and no mechanism for a uniform upper bound has been
proved here.

\subsection{Verification annotations for the spectral formulas}
\label{dev:1981-spectral-checks}

This subsection verifies corrections to established spectral theory. It
does not promote the Wodzicki exposition to a newly open C/D problem.

\begin{proposition}[Correct coefficients in the shifted-zeta identity]\ours
Let \(A\) be an invertible first-order elliptic differential operator on a
closed \(n\)-manifold, with a fixed admissible spectral cut. For a scalar
\(t\) sufficiently small that the cut and invertibility persist, put
\(\delta_\mu=\operatorname*{res}_{s=\mu}\zeta_A(s)\). For \(k\ge0\),
\[
\begin{split}
 \zeta_{A+t}(-k)={}&\sum_{\nu=0}^{k}\binom{k}{\nu}t^\nu\zeta_A(\nu-k)\\
 &+\sum_{\mu=1}^{n}
 \frac{(-1)^\mu\delta_\mu}{\mu\binom{k+\mu}{k}}\,t^{k+\mu}.
\end{split}
\tag{Z1}
\]
The second sum in P001025 (20) omits \(t^\mu\).
\end{proposition}
\begin{proof}
For sufficiently small \(t\) the complex-power binomial expansion,
initially in a trace-class half-plane, is
\[
 \zeta_{A+t}(s)=\sum_{j=0}^\infty
       \frac{(-1)^j(s)_j}{j!}t^j\zeta_A(s+j),
 \qquad (s)_j=s(s+1)\cdots(s+j-1).
\]
One may alternatively take a finite expansion with a remainder of
arbitrarily negative operator order, which is holomorphic near the point
under consideration. For \(0\le j\le k\), the coefficient at
\(s=-k\) is \(\binom{k}{j}\). For \(j=k+\mu\), \(\mu\ge1\),
the Pochhammer product has a simple zero and
\[
 \left.\frac{d}{ds}\frac{(-1)^{k+\mu}(s)_{k+\mu}}{(k+\mu)!}
 \right|_{s=-k}
 =\frac{(-1)^\mu k!(\mu-1)!}{(k+\mu)!}
 =\frac{(-1)^\mu}{\mu\binom{k+\mu}{k}}.
\]
It contributes only when \(\zeta_A(s+j)\) has a pole at \(\mu\).
For a differential operator of order one the positive integer poles are
confined to \(1,\ldots,n\), and the nonpositive integer values are
regular. Multiplying the zero by the pole gives (Z1). The factor \(t^j\)
has remained \(t^{k+\mu}\) throughout. The holomorphic remainder makes
no additional contribution, since its binomial coefficient vanishes at
\(-k\) once the expansion is taken beyond \(k+n\).
\end{proof}

\begin{remark}[Second check and the limit of the inference]\oursComputed
Direct differentiation of the finite Pochhammer product agrees with the
displayed coefficient in all 42 exact rational checks
\(0\le k\le6\), \(1\le\mu\le6\). Setting \(t=0\) independently
detects the notebook's missing powers, whenever the residue combination
is nonzero. Kernel-crossing cases are excluded from the proposition;
adding \(h(A)\) at zero without tracking which zero modes move is not a
substitute for that hypothesis.

The check also prevents an invalid inference about spectral cuts.
Cut independence at \(s=0\) for every shifted operator does not imply
cut independence at every negative integer. In the difference of (Z1)
for two cuts, the case \(k=0\) contains no negative values from which to
start such an induction. The general valid identity is instead
\citet[Theorem 4.2]{ponge2006spectral},
\(\rho(k)=(2\pi i/m)\operatorname{Res}(\Pi P^{-k})\).
Moreover, negative poles of eta and absolute zeta functions occur for
Dirac-type differential operators in odd dimensions
\citep[\S4]{loyamoroianu2012singularities}. For self-adjoint \(P\), the
two cut functions have the form
\(\zeta_+(s)+e^{\pm i\pi s}\zeta_-(s)\); their difference at an
integer is the limit of \(2i\sin(\pi s)\zeta_-(s)\).
A nonzero residue of \(\zeta_-\) therefore gives a nonzero difference,
even though each cut function is regular there. This is why the blanket
claim in P001025(d) is recorded as false without further hypotheses.
\end{remark}

\begin{theorem}[An explicit first-order differential operator with cut-dependent negative value]\ours
\label{dev:1981-negative-cut}
Let \(\mathbb T^3=(\mathbb R/2\pi\mathbb Z)^3\) have its flat metric,
let \(\mu>0\), and let \(\sigma_1,\sigma_2,\sigma_3\) be the Pauli
matrices. On the trivial Hermitian bundle of rank four set
\[
 \alpha_j=\sigma_1\otimes\sigma_j\quad(1\le j\le3),
 \qquad \beta=\sigma_3\otimes I_2,
 \qquad
 P=\sum_{j=1}^3\alpha_j(-i\partial_{x_j})+\mu\beta.
\]
Then \(P\) is an invertible self-adjoint elliptic differential operator
of order one and of Dirac type. Define \(Z_+(s)\) using argument
\(+\pi\) for negative eigenvalues, and \(Z_-(s)\) using argument
\(-\pi\); both branches give argument zero to positive eigenvalues.
The respective branch ranges are \((-\pi/2,3\pi/2)\) and
\((-3\pi/2,\pi/2)\), so the cuts are the negative and positive
imaginary axes. Both cut zeta functions are regular at \(s=-1\), but
\[
 \boxed{Z_+(-1)=-i\pi^2\mu^4,\qquad
        Z_-(-1)=+i\pi^2\mu^4.} \tag{Z2}
\]
In particular, their difference there is \(-2i\pi^2\mu^4\ne0\).
\end{theorem}
\begin{proof}
The four matrices \(\alpha_1,\alpha_2,\alpha_3,\beta\) are
Hermitian, square to the identity, and anticommute in distinct pairs.
Therefore
\[
 P^2=(-\Delta+\mu^2)I_4.
\]
The principal symbol squares to \(|\xi|^2I_4\), proving ellipticity
and the Dirac-type assertion. On the Fourier mode indexed by
\(k\in\mathbb Z^3\), the Hermitian matrix
\(\sum_j\alpha_jk_j+\mu\beta\) has square
\((|k|^2+\mu^2)I_4\) and trace zero. Its eigenvalues are
\(\pm\sqrt{|k|^2+\mu^2}\), each with multiplicity two. Fourier
expansion supplies the self-adjoint realization with domain \(H^1\),
and \(\mu>0\) excludes a kernel. Both stated cuts are admissible for
the operator and its principal symbol.

For \(\operatorname{Re}s>3\), put
\[
 Z_{\mathrm{abs}}(s)=\operatorname{Tr}|P|^{-s}
 =4\sum_{k\in\mathbb Z^3}(|k|^2+\mu^2)^{-s/2}.
\]
The multiplicities and branch conventions give exactly
\[
 Z_\pm(s)=\frac{1+e^{\mp i\pi s}}2\,Z_{\mathrm{abs}}(s). \tag{Z3}
\]
These identities continue meromorphically. To determine the pole of
\(Z_{\mathrm{abs}}\), use its Mellin representation
\[
 Z_{\mathrm{abs}}(s)
 =\frac1{\Gamma(s/2)}\int_0^\infty
 t^{s/2-1}\operatorname{Tr}(e^{-tP^2})\,dt.
\]
Gaussian Poisson summation on this torus yields
\[
\begin{split}
 \operatorname{Tr}(e^{-tP^2})
 &=4e^{-\mu^2t}\sum_{k\in\mathbb Z^3}e^{-t|k|^2}\\
 &=4\pi^{3/2}t^{-3/2}
 \left(1-\mu^2t+\frac{\mu^4t^2}{2}+O(t^3)\right)
   +O\!\left(t^{-3/2}e^{-c/t}\right)
\end{split}
\]
for some \(c>0\) as \(t\downarrow0\). The coefficient of
\(t^{1/2}\) is \(a=2\pi^{3/2}\mu^4\). The corresponding Mellin
integral contributes \(2a/(s+1)\) before division by
\(\Gamma(s/2)\). Since \(\Gamma(-1/2)=-2\sqrt\pi\),
\[
 \operatorname*{res}_{s=-1}Z_{\mathrm{abs}}(s)
 =\frac{2a}{\Gamma(-1/2)}=-2\pi\mu^4. \tag{Z4}
\]
After the displayed heat terms are subtracted, the small-time integral
is holomorphic near \(-1\); the large-time integral is entire because
\(\mu>0\). Thus the pole computation is justified without a formal
exchange of divergent series.

Each prefactor in (Z3) has a simple zero at \(-1\), with
\[
 \left.\frac{d}{ds}\frac{1+e^{\mp i\pi s}}2\right|_{s=-1}
 =\pm\frac{i\pi}{2}.
\]
Multiplying this derivative by the residue in (Z4) proves regularity
and the two values (Z2). The holomorphic part of
\(Z_{\mathrm{abs}}\) contributes zero at that point.
\end{proof}

\begin{remark}[Independent algebraic and normalization checks]\oursComputed
The separate executable \texttt{verify\_negative\_cut.py} checks all
Clifford anticommutators, Hermitian symmetry, the square and the exact
characteristic polynomial, the Mellin residue, and both branch limits.
The heat normalization was independently checked as
\(4\operatorname{vol}(\mathbb T^3)/(4\pi)^{3/2}=4\pi^{3/2}\).
A second mathematical audit reproduced the residue and phase signs.
The symmetric spectrum also gives \(\eta(P;s)=0\) identically:
this example uses the nonzero residue of the absolute zeta function.
It is an explicit verification of the notebook correction, with no
historical priority claim.
\end{remark}

\begin{remark}[The vector-calculus sign check]\oursComputed
For conventional Euclidean operators,
\(\operatorname{div}\operatorname{curl}=0\),
\(\operatorname{curl}\operatorname{grad}=0\), and
\(\operatorname{curl}\operatorname{curl}
=\operatorname{grad}\operatorname{div}-\Delta\).
Direct multiplication gives the corrected square in 1981-558.6 and
1981-559.12. An independent symbolic polynomial calculation replaces
the derivatives by three commuting indeterminates and reproduces both
the erroneous matrix's actual square and the corrected square exactly.
The output is included with the coefficient checks above.
\end{remark}
